\chapter{Conclusions}

The transversity distribution function is one of the leading order PDFs which provides transverse quark polarization in transversely polarized nucleons, which is least constrained from experimental measurements. The polarized proton-proton collision can be accessed via the dihadron channel, in which two unpolarized oppositely charged hadron is detected in the final state. In this channel, the transversity is coupled with the IFF, where the transversity can be described collinearly. To isolate the transversity, the IFF has to be measured independently, for example, in the $e^+e^-$ process. In addition, the other observable required is the unpolarized fragmentation function. Belle experiment has provided a cross section that can provide the quark fragmentation; however, gluon fragmentation has to rely on models resulting in large model-dependent uncertainty in the transversity. Thus the high-precision dihadron azimuthal correlation asymmetry and the unpolarized cross section are the much-needed observables for the model-independent extraction of transversity. 

The $\pi^+\pi^-$ azimuthal correlation asymmetry has been measured using polarized proton-proton collision data at the $\sqrt{s}$ = 200 GeV from STAR run 2015. This dataset provides the most precise asymmetry probing transversity in the valence quark region. The presented $\pi^+\pi^-$ asymmetry measurement corroborates previous STAR measurements, which feature resonance peak at $M_\rho \sim 0.8$ GeV/$c^2$ consistent with the model prediction. A large asymmetry signal in the forward pseudorapidity region, where a high$-x$ quark is probed, suggests the large transversity signal in the valence quark region. A smaller asymmetry signal in the backward region suggests a small transversity signal in the low$-x$ region and purely a valence quark property. 

The unpolarized $\pi^+\pi^-$ cross section differential in \minv in proton-proton collision data at $\sqrt{s}$ = 200 GeV has been measured using STAR run 2012 data for the first time. The measured cross section shows promising  agreement with the PYTHIA cross section and theory prediction. This observable is of much need constraining the gluon fragmentation function. Additionally, the high-precision IFF asymmetries obtained from Run 15, combined with this measurement, will provide a foundation for transversity extraction with greater precision than previously possible.
